\documentclass{anvil-paper}
\usepackage{array}
\usepackage{calc}
\usepackage{fontspec}
\setmainfont{Times New Roman}
\setmonofont{Menlo}[Scale=0.85]
\providecommand{\tightlist}{\setlength{\itemsep}{0pt}\setlength{\parskip}{0pt}}
% Long Lean identifiers in inline code: allow breaks at underscores and dots.
\renewcommand{\texttt}[1]{\path{#1}}
\sloppy

\title{Computational Evidence for a Drop and a Uniform Reduction for Erdős Problem 85}
\author{Claude Fable and GPT Sol \\ Lean Genius project, 2AM Logic}
\date{2026-09-28}

\begin{document}
\maketitle

% Working-manuscript status (kept from the Markdown draft; not part of the paper body):
% **Status.** Working manuscript. Nothing in this document is for external
% distribution before the operator read-through and Zenodo gate. Formal claims
% refer to Lean 4.31.0 with the repository-pinned mathlib. Result A is a
% computational claim resting on a completed two-solver H1 census; it is not a
% proved finite theorem. The conditional finite-drop core is Lean
% checked, but its order-49 nonexistence hypothesis has not been discharged.
% The operator has cancelled certificate replay and production for this paper.
% Theorem B is the Lean-checked conditional reduction to A-REG, which itself
% remains open.

\begin{abstract}
Let \texttt{f(n)\ =\ minDegreeForC4\ n}, the minimum threshold such that every simple graph on \texttt{n} vertices with minimum degree at least \texttt{f(n)} contains a \texttt{C₄}. We report two complementary results. First, the completed computation supports the adjacent-value claim:

\texttt{f(48)\ =\ 8} and \texttt{f(49)\ =\ 7},

which, at the level of computational evidence, is a strict finite drop. The lower side uses a 48-vertex extremal witness checked in Lean through \texttt{native\_decide} (six named \texttt{native\_decide} axioms in total for the two witness endpoints; the literal \texttt{\#print\ axioms} output from a cold rebuild is banked as \texttt{AXIOM\_AUDIT\_COLD\_20260927/}); an independently computed isomorphism check finds it non-isomorphic to the previously recorded Afzaly--McKay witness. For the order-49 upper side, structural reductions and archived SAT evidence have different verification levels; every one of the 1,161 residual H1 instances returned UNSAT from two independent solvers without proof certificates (1,160 as whole instances, one as an exact partition into 36 cubes). Second, we reduce a negative answer to Erdős Problem 85 to one uniform graph-theoretic proposition, \texttt{BinarySquareRegularExclusion} (A-REG). Lean verifies the entire implication

\texttt{BinarySquareRegularExclusion\ →\ ¬\ Erdos85Question}.

A-REG is a live hypothesis, not a proved theorem or a forecast. In particular, the exact Lean results at orders 15 and 16 exhibit no drop and refute the \texttt{q\ =\ 4} analogue, while the published Boza table reports the same no-drop behavior at orders 35 and 36. We therefore present the operator's plane-order interpretation as a genuine rival hypothesis. The extensive negative map records which determinant, spectrum, packing, incidence, and finite-census routes fail to settle A-REG, isolating the connected and mixed non-bipartite residues that remain open.
\end{abstract}

\section{Main results}\label{main-results}

\subsection{Result A --- computational evidence for a 48-to-49 drop}\label{result-a-computational-evidence-for-a-48-to-49-drop}

The combined evidence supports the adjacent-value claim

\texttt{minDegreeForC4\ 48\ =\ 8\ ∧\ minDegreeForC4\ 49\ =\ 7},

and consequently the proposed strict drop

\texttt{minDegreeForC4\ 49\ \textless{}\ minDegreeForC4\ 48}.

This is a computational result, not an unconditional Lean theorem. The census rule was that any row reaching its declared solver cap would be printed as open; after escalating caps and one cube partition, no row is open. Even with every row UNSAT, proof certificates would still be needed for a formal theorem.

The already checked finite-drop core is \texttt{minDegreeForC4\_fortyEight\_fortyNine\_exact\_checked} and \texttt{minDegreeForC4\_fortyNine\_lt\_fortyEight\_checked}; both take the order-49 nonexistence hypothesis. The conditional generated endpoints are specified, under the planned names \texttt{minDegreeForC4\_fortyEight\_fortyNine\_exact\_of\_generatedSevenBaseCertificates} and \texttt{minDegreeForC4\_fortyNine\_lt\_fortyEight\_of\_generatedSevenBaseCertificates}, in a generated module that is not yet present in the source tree. They would prove the finite drop only after every input is supplied, the generated module lands, and a cold build and literal axiom audit pass. No conjectural uniform hypothesis enters this finite computation. Separately, an independently computed NetworkX isomorphism comparison of archived graph6 artifacts finds that the checked order-48 extremal graph is not isomorphic to the Afzaly--McKay record; this novelty check is not part of the Lean-elaborated theorem, and that theorem's witness side rests on \texttt{native\_decide}, not on a standard-axiom-only kernel proof.

\subsection{Evidence levels and formal obligations for Result A}\label{evidence-levels-and-formal-obligations-for-result-a}

The evidence classes below are not a partition: the historical 96 H1 roots, the existing certificate-object inventory, and the Phase B root set use different selections of the same capacity tags. No row is promoted from a solver verdict or an archived proof object to a kernel-checked theorem.

{\def\LTcaptype{none} % do not increment counter
\begin{table}[htbp]\centering\small
\begin{tabular}{@{}
  >{\raggedright\arraybackslash}p{(\linewidth - 4\tabcolsep) * \real{0.3333}}
  >{\raggedright\arraybackslash}p{(\linewidth - 4\tabcolsep) * \real{0.3333}}
  >{\raggedright\arraybackslash}p{(\linewidth - 4\tabcolsep) * \real{0.3333}}@{}}
\toprule
\begin{minipage}[b]{\linewidth}\raggedright
Class
\end{minipage} & \begin{minipage}[b]{\linewidth}\raggedright
Current evidence
\end{minipage} & \begin{minipage}[b]{\linewidth}\raggedright
Limit for this paper
\end{minipage} \\
\midrule
\bottomrule
Order 48 and 49 lower witnesses & Finite graphs in \texttt{Erdos85FiniteDropWitnesses.lean}, elaborated in Lean via \texttt{native\_decide}; the cold-rebuild \texttt{\#print\ axioms} audit (\texttt{AXIOM\_AUDIT\_COLD\_20260927/}) lists three \texttt{native\_decide} axioms per witness endpoint, and independent edge-list audits (168 edges, all pair codegrees ≤ 1) corroborate both graphs & Establish the lower sides only; not standard-axiom-only. \\
H3/H5 & Reviewed paper/computation cover and archived local proof receipts & Generated formal aggregate and public axiom audit are absent. \\
H7 & All 28 surviving structural roots covered after 15 singleton-capacity exclusions (\texttt{H7\_CLOSURE\_20260915.md}) & Paper/computation closure; the Lean evidence-vector arguments remain uninstantiated. \\
H1 historical overlay & 96 reviewed historical cases with archived \texttt{drat-trim} verification (\texttt{phase\_b\_historical\_overlay\_96/}) & Historical checks are not a current kernel replay. \\
H1 existing bank & 12,019 historical ready certificate inputs in the replay sizing model & Object and metadata availability do not establish a completed kernel replay. \\
H1 residual roots & 1,161 Phase B roots without a listed certificate object: 1,160 returned UNSAT from Kissat 4.0.4 and then CaDiCaL 3.0.1 under declared caps (4, 12 or 24 hours per solver over four passes), and one, \texttt{h1\_81494a6ef36d3ec9}, which defeated every whole-instance cap, returned UNSAT on all 36 leaves of an exact cube partition under both solvers; receipts re-derived by the reviewed summarizer and cube checker (\texttt{phase\_b\_h1\_census\_20260927/}) & Verdict-only evidence: two solvers agreeing under a cap is evidence about solver behaviour, not a proof; no certificate is produced. \\
H1 capacity-grid slots & The 1,288 gap slots of the banked capacity snapshot are an overlapping decomposition: 1,158 are residual roots above, 96 are the historical overlay, 34 lie outside the frozen Phase B source and were solved through their own reviewed route (\texttt{phase\_b\_h1\_gap130/}, all 34 UNSAT under both solvers, Kissat 0.3 to 1.7 hours each). The reviewed gap auditor (\texttt{phase\_b\_h1\_census\_20260927/h1-gap1288-audit.json}) reports 1,191 slots with fresh two-solver UNSAT verdicts, 96 with inherited 2026-08 certificate verification, and 1 slot, the cube-partitioned row, which it cannot read and lists as UNKNOWN & By the auditor's own definition only fresh verdicts count, so it prints 97 ``open'' slots: the 96 inherited rows and the cube row. Every slot has evidence; none has a certificate produced by this paper. \\
\end{tabular}
\end{table}
}

The residual-root count is final and comes from exact tag and CNF joins of the solver receipts by the reviewed summarizer, not from a sum of overlapping inventory counts. The complete set of two-solver UNSAT verdicts supports Result A computationally; it leaves the formal nonexistence theorem unproved.

The mathematical dependency chain is independent of how the computation is scheduled. The checked witness module \texttt{Erdos85FiniteDropWitnesses.lean} supplies \texttt{boza48\_degreeSeven\_witness} and \texttt{orderFortyNine\_degreeSix\_witness}. Its exact-value theorem still takes \texttt{hno49\ :\ ¬\ C4FreeMinDegreeWitness\ 49\ 7}. No completed-certificate count by itself supplies this hypothesis.

The source-level consumer \texttt{not\_c4FreeMinDegreeWitness\_fortyNine\_seven\_of\_smallHighLratChecks} in \texttt{Erdos85OrderFortyNineSmallHighVerifiedFrontier.lean} requires four inputs:

{\def\LTcaptype{none} % do not increment counter
\begin{table}[htbp]\centering\small
\begin{tabular}{@{}
  >{\raggedright\arraybackslash}p{(\linewidth - 2\tabcolsep) * \real{0.5000}}
  >{\raggedright\arraybackslash}p{(\linewidth - 2\tabcolsep) * \real{0.5000}}@{}}
\toprule
\begin{minipage}[b]{\linewidth}\raggedright
Sector
\end{minipage} & \begin{minipage}[b]{\linewidth}\raggedright
Required input to this consumer
\end{minipage} \\
\midrule
\bottomrule
H1 & \texttt{OrderFortyNineStratumExcluded\ 1} \\
H3 & Checked LRAT proofs for both representative indices \texttt{index\ ≤\ 1} \\
H5 & Checked LRAT proofs for all three representative indices \texttt{index\ ≤\ 2} \\
H7 & \texttt{OrderFortyNineStratumExcluded\ 7} \\
\end{tabular}
\end{table}
}

For H3 and H5, the checked formulas are precisely \texttt{orderFortyNineGeneratedCanonicalSatCnf} applied to \texttt{threeHighRepresentativeMasks} or \texttt{fiveHighRepresentativeMasks}. This is an interface description, not a claim that five monolithic certificates are available or that five jobs suffice. The finer cube route instead supplies seven base-CNF \texttt{Unsat} proofs to \texttt{not\_c4FreeMinDegreeWitness\_fortyNine\_seven\_of\_smallHighCubeBaseUnsat} in \texttt{Erdos85OrderFortyNineSmallHighCubeGridTerminal.lean}. It reaches the same \texttt{hno49} conclusion without constructing the five existential LRAT arrays of the older consumer. The finite-drop core accepts either route. H1 capacity rows still need the full row-to-stratum assembly. The H7 structural ledger \texttt{H7\_CLOSURE\_20260915.md} covers all 28 roots that survive the 15 singleton capacity exclusions at the paper/computation level. This does not yet instantiate the Lean capstone \texttt{orderFortyNineStratumExcluded\_seven\_of\_emptyCubeEvidenceVectors}, whose statement still takes four evidence-vector arguments of lengths 19/15/7/2. H1 still has a root-coverage campaign; H7's formal evidence bridge and the final generated aggregate remain open proof obligations.

The companion closure inventory records the chosen route for each sector, its missing artifacts, and replay obligations. Its counts and cost estimates are operational evidence, not additional theorem hypotheses. A future formal closure would require the exact cover and semantic bridges, certificate replay for every necessary row, the generated aggregate, and a cold build with literal axiom audits of the two public exact-value/drop endpoints. None of these steps follows from verdict-only UNSAT results.

\subsubsection{Where the residual trust sits}\label{where-the-residual-trust-sits}

The weakest link in any SAT-based nonexistence result is the encoding, not the solver. Three facts bound that risk here. The H1 instance generator is itself a Lean program, \texttt{Proofs/Erdos85OneHighV2CnfEmit.lean}, compiled to the native \texttt{v2cnf} binary (sha256 \texttt{4bd9604c…}) that every run pins; the emitted DIMACS is regenerated and compared by the binary's own check mode before any solver starts; and in 1,322 of the 1,412 input preparations of the census the emitted CNF matched, byte for byte, the hash recorded by the independent 2026-08 producer. What the encoding does not yet have is a Lean-elaborated semantic bridge: for H3 and H5 the checked formulas are the Lean terms \texttt{orderFortyNineGeneratedCanonicalSatCnf} applied to the representative masks, whereas for H1 the row-to-stratum assembly is an open formal obligation. An encoding error would be systematic and invisible to solving; the h305 case study below records one such error that review caught. After the encoding come the solvers themselves (two independent UNSAT verdicts are evidence about solver behaviour, not a proof) and the uninstantiated H7 Lean capstone.

\subsection{Cost to verify Result A formally}\label{cost-to-verify-result-a-formally}

The cancelled replay plan is retained as a reproducible cost model, not an active launch plan. For 12,019 historical ready certificate inputs, its byte-weighted model estimates 4,831 compile box-hours. With the plan's 25\% noncompile allowance, 10\% spot-loss factor, and eight bootstrap hours, this becomes about 6,651 box-hours: \textbf{17.32 ideal allocated days on 16 hosts} and \textbf{about \$1,276 of infrastructure} at the sampled spot, gp3 and IPv4 rates (\texttt{sat49/H1\_REPLAY\_SPOT\_16\_BUDGET\_PLAN\_20260916.md}). Retrieval is additional: a full 6.06 TB pass at the cited Glacier Instant Retrieval rate of \$0.03/GB would cost about \$182 before requests, repeat reads, output storage or other services. No replay wave is authorized by this estimate.

The missing-certificate production cost is unresolved. An early \textbf{rough \$2,600 proxy} multiplied 1,288 gaps by a 10-hour claim-to-ledger interval; the worker logs show that interval includes queue and pipeline time, so it must not be used as a solver-time forecast. The later v3 readback finds 574 verified-UNSAT solves with a 4,011-second median and a 4,539-second mean, and 22 UNKNOWN rows at a 14,400-second cap (\texttt{sat49/H1\_V3\_SOLVER\_TIMING\_20260916.md}). The unproved rows are selected for difficulty, and proof logging, trimming, splits and retries have no measured cost model here. Thus the gap-certificate figure is a historical planning proxy, not a credible fixed price. A public requester-pays copy of the existing certificate bank could let others fund independent verification, subject to the operator's publication gate and an exact release manifest.

\subsubsection{A cheaper route: cube-partitioned certificates (projection)}\label{a-cheaper-route-cube-partitioned-certificates-projection}

The price above assumes one certificate per instance, checked whole. The archived bank shows why that is expensive: across 12,102 verified H1 rows the compact LRAT proof grows at a nearly constant 5 to 6 GB per Kissat solve-hour (median 1.2 GB, 90th percentile 4.0 GB, largest 26.5 GB, 22.6 TB in total), and the replay pilot checked about half a megabyte of gzipped certificate per second, so a kernel check costs roughly as long as the solve that produced it. The largest files force one checker per 64 GB host, which is what turns 4,800 compile hours into a month of fleet time. The 1,161 residual roots settled by the verdict-only census are heavier still: 1,159 UNSAT verdicts needed 2,695 Kissat core-hours (median 1.9 hours per row; 120 rows above 4 hours, six above 12 hours), so whole-instance certificates for them would total about 15 TB, and the hardest rows would produce single files of 60 to 130 GB that no checker can hold.

Cube-and-conquer changes the shape of that cost. Fixing the values of a few variables splits an instance into cubes that partition its search space; the instance is UNSAT exactly when every cube is. Each cube yields its own small certificate, so checkers need a few gigabytes rather than 64, eight or more checks share one host, a lost spot instance costs one cube rather than a 26 GB job, and a short Lean lemma (a complete tree of UNSAT cubes refutes the base formula) composes the leaves. The trivial fraction is large: the hardest residual row, which failed every whole-instance cap up to 24 hours, was refuted by an adaptive cube tree of 36 leaves (35 splits, depth at most 8) in 9.7 Kissat and 6.2 CaDiCaL core-hours, under five hours of wall clock on 24 cores; 29 of the 36 leaves fell within 15 minutes and the hardest needed 52 minutes. Proof size still tracks solver work, so the total byte count is the open question; splitting can lower it where a single CDCL run thrashes and raise it where work is duplicated across cubes.

Under the measured rates, a cube-certified census of the 1,161 residual roots would need about 6,700 core-hours of proof-logged solving and trimming (roughly \$70 to \$100 on spot), about 360 host-hours of kernel checking packed eight to a host (roughly \$70), and 4 TB of gzipped certificates (about \$16 per month in cold storage), against \$2,000 to \$2,500 for the whole-instance plan. These are projections from the archived bank and one live cube tree, not receipts; the Lean-side composition lemma and per-cube proof emission are not yet built. We record the route because it makes the price of certainty an order of magnitude more approachable for whoever takes Result A to a theorem.

\subsection{Theorem B --- a one-proposition reduction of the infinite problem}\label{theorem-b-a-one-proposition-reduction-of-the-infinite-problem}

Define \texttt{BinarySquareRegularExclusion} to assert that, for every \texttt{k\ ≥\ 3}, no simple \texttt{C₄}-free \texttt{2\^{}k}-regular graph exists on \texttt{4\^{}k} vertices. Lean proves

\texttt{not\_erdos85Question\_of\_binarySquareRegularExclusion\ :\ \ \ BinarySquareRegularExclusion\ →\ ¬\ Erdos85Question}.

The existence jaw comes from deleting the absolute nucleus of the even-characteristic polarity graph; tight square-order candidates are forced to be regular; and both unit defect components and bipartite defect components are excluded. The remaining honest residue is A-REG-NONBIP. Its one-component socket is stated by \texttt{BinarySquareConnectedNonbipartiteExclusion}, but no uniform theorem currently routes every disconnected or mixed non-bipartite partition into that socket. Thus Theorem B is a complete checked reduction, not a proof of A-REG.

\subsection{Contributions and collaboration}\label{contributions-and-collaboration}

Claude Fable led the certification pipeline and cold-integration verification, the existence halves of the census, fleet operations, the September 2026 verdict-only census with its cube-partition tooling, §8, and the final scope review. GPT Sol developed the structural reductions, negative map, and independent audits. The human operator set compute policy, research priorities, authorship, scope, and the final read-through gate. The shared room infrastructure supplied durable coordination, claims, review gates, and a reproducible transcript; it is acknowledged as infrastructure, not as an author. High-variance proposals were admitted only after independent Lean elaboration or exact certificate replay.

\begin{center}\rule{0.5\linewidth}{0.5pt}\end{center}

\section{0. Mathematical status: one uniform axiom remains}\label{mathematical-status-one-uniform-axiom-remains}

Let \texttt{f(n)\ =\ minDegreeForC4\ n}, the minimum threshold such that every simple graph on \texttt{n} vertices with minimum degree at least \texttt{f(n)} contains a \texttt{C₄}. Erdős Problem 85 asks whether \texttt{f} is eventually nondecreasing. The formal reduction is complete: the negation is equivalent to the existence of arbitrarily late strict drops (\texttt{erdos85Negation\_iff\_not\_question}), and a \texttt{q}-regular witness on \texttt{q²−1} vertices together with nonexistence at \texttt{(q²,q)} produces such a drop (\texttt{PlaneOrderDropWitness.strict\_drop}). Cofinally many such pincers therefore refute the problem (\texttt{not\_erdos85Question\_of\_cofinalPlaneOrderDropFamily}).

For \texttt{q\ =\ 2\^{}k}, \texttt{k\ ≥\ 3}, the existence jaw is already uniform: delete the absolute nucleus from the even-characteristic polarity graph (\texttt{Polarity.c4FreeMinDegreeWitness\_even\_delete\_absolute\_nucleus}), and the resulting orders are cofinal (\texttt{cofinalEvenFieldSquareExclusion\_of\_binary}). On the nonexistence jaw, the normalized square-order reduction is an equivalence (\texttt{squareOrderTightCoreExists\_iff\_witness}, \texttt{binarySquareOrderTightCoreExclusion\_iff}), while parity forces every such tight core to be regular (\texttt{squareOrder\_regular\_of\_even}). Consequently the whole infinite theorem follows from one proposition:

\begin{quote}
\textbf{A-REG (\texttt{BinarySquareRegularExclusion}).} For every \texttt{k\ ≥\ 3}, there is no simple \texttt{C₄}-free \texttt{2\^{}k}-regular graph on \texttt{2\^{}k\ ·\ 2\^{}k\ =\ 4\^{}k} vertices.
\end{quote}

This is not shorthand for several hidden assumptions. Lean proves directly that A-REG implies the normalized tight-core exclusion (\texttt{binarySquareOrderTightCoreExclusion\_of\_regularExclusion}) and hence the negation of Erdős 85 (\texttt{not\_erdos85Question\_of\_binarySquareRegularExclusion}). Thus every arrow from A-REG to the headline theorem is cold-green; A-REG itself remains an axiom/conjecture, not a proved result (outline v2.64, §A.5; room msg 31964).

The defect operator makes the residue precise. For a hypothetical regular graph with adjacency matrix \texttt{A},

\texttt{A²\ =\ (q−1)I\ +\ J\ −\ D},

and \texttt{A} commutes with \texttt{D} (\texttt{adjMatrix\_sq\_eq\_sub\_secondOrderDefect\_of\_regular}, \texttt{adjMatrix\_comm\_secondOrderDefect\_of\_regular}). Its defect components have orders \texttt{q\ m\_c}, with \texttt{Σ\_c\ m\_c\ =\ q} (\texttt{binarySquare\_regular\_exists\_defectComponent\_partition}), and every vertex of component \texttt{c} has exactly \texttt{m\_c} graph-neighbours within that component (\texttt{binarySquare\_regular\_degree\_induce\_defectComponent\_eq\_part}). Unit parts are impossible for even \texttt{q} (\texttt{binarySquare\_regular\_no\_sizeQ\_defectComponent\_of\_even}), as are bipartite defect components when \texttt{4\ \textbar{}\ q} (\texttt{binarySquare\_regular\_no\_bipartite\_defectComponent}). What remains is therefore exactly \textbf{A-REG-NONBIP}: all partitions \texttt{q\ =\ Σ\ m\_c} with \texttt{m\_c\ ≥\ 2} and every defect component non-bipartite.

The one-part subcase has an especially sharp equivalent formulation:

\begin{quote}
\textbf{NONBIP-CONNECTED.} For binary \texttt{q\ ≥\ 8}, every loopless \texttt{q}-regular \texttt{C₄}-free adjacency matrix \texttt{A} of order \texttt{q²} is singular. Equivalently, its defect graph \texttt{D} is not connected, since \texttt{dim\ ker(A)\ =\ numberOfComponents(D)\ −\ 1} (\texttt{binarySquare\_regular\_finrank\_adj\_kernel\_eq\_card\_components\_sub\_one}).
\end{quote}

The formal branch socket is \texttt{BinarySquareConnectedNonbipartiteExclusion}. Its nonbipartite hypothesis is now discharged literally for every regular binary-square candidate by \texttt{binarySquare\_twoPow\_regular\_defectGraph\_not\_bipartite} (with the stronger per-component coloring exclusion \texttt{binarySquare\_twoPow\_regular\_no\_bipartite\_defectComponent}). Its connected hypothesis is different: there is currently no uniform Lean theorem routing all disconnected or articulated defect graphs into this socket. Such a routing is closed only in the separate \texttt{q\ =\ 9} B.3 analysis. Thus the connected socket and the mixed nonbipartite partitions are sibling open cases for binary \texttt{q}; connectivity must not be cited as a proved uniform reduction.

This formulation includes all graph hypotheses; proving singularity from generic regularity or spectrum alone would not suffice. The campaign's current mathematical frontier is the implication from the full symmetric \texttt{C₄}-free incidence completion to that singularity statement. Mixed non-bipartite partitions remain a sibling subcase, not a consequence of the connected case (outline v2.64, §E).

\subsection{Negative map: discarded routes are part of the result}\label{negative-map-discarded-routes-are-part-of-the-result}

The search did not merely fail to finish several familiar approaches; it produced exact countermodels or reductions showing why they do not close the remaining node.

\begin{itemize}
\tightlist
\item
  The identity \texttt{det(A)²\ =\ q⁴\ τ(D)} makes a square spanning-tree count necessary when \texttt{D} is connected. Controls show that this condition alone does not force singularity (outline §A.5.3(i), \texttt{0ed91c72d6}). This does not exclude stronger integral or local arithmetic tests: later examples fail precisely such tests.
\item
  Real spectral controls satisfy the particular scalar conditions recorded in their reports (\texttt{f6ee2ed421}; divergence \#66). They do not establish integral adjacency realizability or compatibility with every odd-trace condition. Later Hoffman-diagonal and odd-power arguments reject particular proposed spectra; these rejections and the remaining scope are recorded in outline versions 2.66--2.69.3.
\item
  Generic packing, Hall, code-LP and fractional-transversal statements omit the completion of all rows to one symmetric adjacency matrix. The strongest faithful self-polar partial interface already has an exact half-integral survivor at \texttt{q=6}, so integrality does not follow from those hypotheses (\texttt{e401f3034f}; room msg 31951).
\item
  The proposed P8 inverse-sign separation is impossible: exact product-sign identities force both signs among the relevant inverse entries. The false terminal was retracted, not weakened (\texttt{19f227dced}; outline v2.62, room msgs 31962--31963).
\item
  Pfaffian/mod-2 valuation, Ihara, exterior-power and Lefschetz variants reduce to the same determinant or spectral data. Coherent-configuration, projective-completion and line-graph thresholds lack the required bridge; the latter premise already fails the \texttt{q=4} incidence control and connected spectral controls. The ten-route audit and its no-survivor verdict are banked in \texttt{NONBIP\_CONNECTED\_LITERATURE\_DIVERGENCE.md} at \texttt{cf2243a6e8} (room msgs 31957 and 31962).
\end{itemize}

The subsequent cycle produced several further boundaries with independently reviewed scope (\texttt{CUTS\_LEDGER\_DRAFT.md}, rows 172--186):

\begin{itemize}
\item
  A core-Lean partial common-neighbour reconstruction identifies the exact graph problem in operation form. A separate prose construction, checked at finite field examples, refutes same-carrier Hamming stability based only on the total central law's failure rate. It does not refute every possible additional operation identity (\texttt{PARTIAL\_CENTRAL\_OPERATION\_AUDIT.md}).
\item
  The interval circulant defect family with steps \texttt{\{q²/2,\ ±1,\ …,\ ±(q/2−1)\}} is excluded for every binary \texttt{q\ ≥\ 4} by a reviewed cyclotomic argument (\texttt{l6-cyclic-trace/UNIFORM.md}). Specific order-256 defects also fail rational congruence to the identity form, as witnessed by local Hasse invariants (\texttt{l6-representability/}). The recorded D16 additionally fails the mod-2 alternating-root test; the interval \texttt{r\ =\ 4} example passes the 2-adic representability test but fails at odd primes. These are family exclusions, with no reduction of arbitrary defect graphs to those families.
\item
  No proper divisible design graph has odd degree \texttt{k\ ≥\ 3} and common-neighbour parameters \texttt{λ₁\ =\ 0,\ λ₂\ =\ 1}. Reviewed geometric arguments also rule out loopless polarities preserving the Baer-deleted incidence structure. Matching-repair and retained-pencil bounds rule out stated support and deletion budgets. None excludes arbitrary global repairs (\texttt{baer-repair-gate/}, \texttt{BAER\_POLARITY\_EXTENSION.md}, \texttt{retained-subplane-gate/}).
\item
  Explicit regular connected nonbipartite defect graphs pass the full alternating symmetric square-root test over \texttt{F₂}, but the same family fails the determinant-square condition uniformly. This separates the two necessary tests; it supplies no survivor of their conjunction (\texttt{mod2-root-gate/}).
\item
  The stated five-vertex flag relaxation is feasible uniformly for binary \texttt{q\ ≥\ 16} after repairing the original witnesses to satisfy \texttt{tr(A\ D²)\ ≥\ 0}; this includes the listed 23 identities and eight Gram families, not every possible five-vertex constraint (ledger row 175). The aggregate mod-3 census is feasible for every odd exponent \texttt{k\ ≥\ 5} (row 174), and the projected fixed-shore bound \texttt{z\ ≤\ q−1} is refuted by a uniform \texttt{z\ =\ q} construction (row 176). The specified enlargement of an even-order polarity graph to degree \texttt{q+1} on \texttt{(q+1)²−3} vertices requires at least \texttt{q²/4} old-edge deletions for \texttt{q\ ≥\ 16} (row 180); it excludes repairs with only linearly many deletions, not all repairs.
\end{itemize}

The algebraic and geometric arguments in these reports are reviewed prose unless a named Lean result is explicitly identified. Their computational checks are evidence for their stated examples, not universal certificates.

These are scope statements, not impossibility theorems about every future variant. They say exactly which advertised hypotheses were too weak and identify the missing input: a new invariant that uses the simultaneous completion of \emph{all} rows to a symmetric \texttt{q}-regular \texttt{C₄}-free \texttt{A}. Recording that boundary prevents a finite census, a partial incidence model, or a conditional terminal from being mistaken for the uniform theorem.

\begin{center}\rule{0.5\linewidth}{0.5pt}\end{center}

\section{1. Verification asymmetry}\label{verification-asymmetry}

The campaign separates the cost of proposing a claim from the cost of admitting it to the record. Proposals may be fast and speculative; admission requires either source elaboration under the pinned Lean toolchain or replay of a checked certificate against the exact formal CNF. The distinction is visible in the inverse-potential episode. At 17:46, a root-summed identity was announced as \texttt{tr(A⁻¹)=q} (room msg 31890). Five minutes later its author re-expanded the all-pairs term, found that it was \texttt{q·1ᵀA⁻¹1=q²} rather than \texttt{q·tr(A⁻¹)}, and withdrew the consequence (31895). No Lean wrapper or outline node had been built, so the correction cost minutes and left no false theorem behind.

The same asymmetry governed the order-64 \texttt{h305} endpoint. The first certificate family used 80 owners, but the graph semantics required 88: the missing eight were the antipodal shore pairs. The mismatch was detected by comparing the formal shore modes with the certificate universe, before the endpoint was called proved (room msgs 31664, 31674, 31697; outline v2.64, §A.5.2). The corrected chain ends in \texttt{false\_of\_h305\_source\_or\_transported} and \texttt{muNegThreeZeroFiveEndpointCallback\_false}; the integrator then rebuilt the full chain cold and printed the exact six Owner88 certificate axioms (31809, 31845, 31882). High-variance search is safe here because a plausible certificate is evidence only after semantic identity and kernel replay.

\section{2. Adversarial diversity}\label{adversarial-diversity}

Different agents are useful only when they test different failure modes. Review \#939 did not merely repeat the inverse-potential derivation: it re-derived P1--P7 and then inspected the exceptional root term. That audit found that the proposed sign dichotomy could be satisfied by the root itself and forced the corrected domain \texttt{V\ \textbackslash{}\ (N(y)\ ∪\ \{y\})} (31868--31874). The author then attacked the corrected statement and proved that its P8 sign separation was itself impossible: a suitable defect-neighbour block has zero potential sum, so it cannot be strictly negative pointwise. A second agent caught an overbroad sentence in that retraction---defect edges may also be triangle-free graph edges---and supplied the necessary choice from \texttt{N\_D(y)\ \textbackslash{}\ N\_A(y)} (31905, 31907, 31913). The final correction is banked in \texttt{19f227dced} and \texttt{0ae7069f40}.

Independent checking also corrected campaign accounting. A manuscript audit initially described the thirteen order-49 inputs as five one-high cells plus one seven-high cell. The host manifest listed the actual inputs: four H3 scouts, three H5 cells, and six H7-t0 cubes. The author retracted that sentence within two minutes and corrected the draft (31989, 31994, 31997). The lesson is narrower than ``use multiple models'': the reviewer must inspect an independent invariant---the exact host manifest, a semantic universe, or a boundary term---rather than replay the author's narrative.

\section{3. The structure--compute exchange rate}\label{the-structurecompute-exchange-rate}

The order-49 campaign gives a current, measured exchange. Structural normalization reduces seven H3/H5 monoliths to two checked cover formulas and a \texttt{7×8} grid per cell. Lean proves the accounting---392 positive cubes plus fourteen covers---in \texttt{orderFortyNineSmallHigh\_positiveCube\_job\_count}, and proves their exhaustive composition in \texttt{orderFortyNineSmallHigh\_unsat\_of\_checkedCubeGrid}. The host therefore runs 406 bounded jobs instead of trusting a solver-side cube list (31994). When the campaign was fired, a separate audit found that these grid results did not fit the older five-monolithic-LRAT socket: the cube stack used the VariableHigh CNFs and returned \texttt{CNF.Unsat}, while the final socket expected canonical \texttt{LRAT.check}s. The missing semantic path was formalized as \texttt{not\_c4FreeMinDegreeWitness\_fortyNine\_seven\_of\_smallHighCubeBaseUnsat} (\texttt{d0a17f358b}; room msg 32032) before the first verdict was needed.

The \texttt{h305} repair shows the reverse exchange. Six honest 88-owner CNFs were all externally UNSAT in seconds, but the expensive work was not search; it was proving that the graph semantics map to those exact literals and modes. That chain runs through \texttt{h305\_crossExteriorSplit\_of\_profile}, \texttt{muNegThreeZeroFiveCorrectFiniteSemantics\_false}, and \texttt{false\_of\_h305\_source\_or\_transported} (outline v2.64, §A.5.2; 31882). Compute locates a finite obstruction; structure determines whether the obstruction says anything about the graph.

\section{4. Persistence and the shape of a long campaign}\label{persistence-and-the-shape-of-a-long-campaign}

The campaign state lives in four independent records: Lean sources and git commits, the room transcript, the proof outline, and durable certificate storage. The outline is not a retrospective summary. Operator goal \#25 made it the allocation instrument: every root-to-leaf node is labelled PROVEN, CERT, AXIOM, or GAP, and workers choose lanes against that tree (outline v2.64, §G; goal \#25). This prevented a long finite endpoint from silently becoming the whole project: order-64 \texttt{h305} was parked when its premise was underived, then reopened by goal \#38 after the operator chose the endpoint audit. The reopened audit discovered the 80/88 mismatch, rebuilt the honest universe, and closed the endpoint in one day at \texttt{802f6d79d3} and \texttt{778a2e1595} (31664--31682, 31849). Persistence here means preserving enough state to restart from the actual boundary, not merely keeping a process alive.

The same rule governs compute. Goal \#39 required a pre-fire manifest with all thirteen input hashes, tool hashes, job order, caps, storage exclusions, and deletion policy. The posted manifest verified all thirteen SHA-256 entries before launching a solver (31994). A verdict without that durable provenance is not campaign progress, even if the process printed \texttt{UNSAT}.

\section{5. Honest scoping}\label{honest-scoping}

Every status word names its scope. The callback theorem \texttt{muNegThreeZeroFiveEndpointCallback\_false} carries the foundational axioms plus exactly six disclosed Owner88 \texttt{native\_decide} checks. The broader \texttt{orderSixtyFour\_regular\_sizeTwo\_signedJoint\_false\_of\_connected} additionally inherits the older FullClosure certificate family; the integrator printed both lists separately rather than reporting the smaller list for the larger theorem (31845, 31846, 31882). It closes the disconnected size-two subtree at order 64, not NONBIP-CONNECTED and not A-REG.

The same discipline keeps Result A at its computational evidence level. The operator cancelled certificate production and replay for this paper; the two-solver H1 census was the remaining empirical gate and is complete. A completed census does not prove eventual monotonicity false, since one drop does not settle an eventual property. This manuscript therefore distinguishes a computationally supported 48-to-49 drop from the Lean-checked conditional reduction A-REG-to-\texttt{¬\ Erdos85Question}.

\section{6. The thin human role}\label{the-thin-human-role}

The human role is thin but decisive at branch points. Goal \#36 did not suggest a lemma; it imposed a stuck test: if workers cannot name every link from recent banks to the terminal, they must stop, search outside first, diverge widely, and report refutations as results. That rule ended the inverse-potential lane after P8, fractional-cover integrality, and divergence \#68 all failed (31951--31964). Goal \#38 separately chose to reopen the parked \texttt{h305} endpoint; goal \#39 authorized the certificate campaign that closes the currently launched H3/H5/H7 finite cells but cannot decide 48/49 without a separate H1 closure; goal \#40 commissioned this manuscript in parallel. These are portfolio and governance decisions. The mathematical formal claims still pass through Lean or certificate replay, and nothing goes external before operator review (31965, 31970).

\section{7. Why Erdős problems}\label{why-erdux151s-problems}

Erdős 85 exposes several proof currencies at once: an elementary extremal statement, uniform finite-field witnesses, exact graph reductions, spectral and incidence structure, and finite certificate endpoints. Each partial claim has a precise interface. A-REG would close the uniform theorem through \texttt{not\_erdos85Question\_of\_binarySquareRegularExclusion}; the order-49 campaign would close one finite drop through the small-high socket; and the order-64 work can honestly report a closed signed-joint subtree while leaving the connected defect frontier open (outline v2.64, §0, §A.5, §B.1). That granularity makes the problem a useful test of whether a machine collaboration can accumulate durable mathematics without confusing a large amount of verified local work with a solved root problem.

\begin{center}\rule{0.5\linewidth}{0.5pt}\end{center}

\section{Silence is not success (a recurring failure class)}\label{silence-is-not-success-a-recurring-failure-class}

Silence fails at the mathematical level too. The false trace identity in 31890 looked plausible because no term visibly objected; an explicit re-expansion produced the positive evidence of cancellation and triggered the retraction at 31895. P8 survived a thousand sampled \texttt{q=4} models only because those models were singular and therefore never instantiated the inverse hypothesis; the correct response was \texttt{UNKNOWN}, not support (31886). A direct block-sum argument then proved that P8 was impossible in every hypothetical survivor (31905--31913).

The same standard applies operationally. A solver cell counts only when its ledger line contains a verdict, return code, input hash, proof-check status, LRAT hash, and upload status; a missing line is not an UNSAT. A Lean file counts only after the named target elaborates and its \texttt{\#print\ axioms} output is inspected. On 25 August the \texttt{h305} dependency chain ran for nearly an hour, but the room did not infer success from the absence of errors: three agents waited for the final file check and separately reported its axiom scope (31845, 31846, 31882). Positive artifacts, rather than quiet logs, are the unit of trust.

\section{Methods (summary)}\label{methods-summary}

\begin{itemize}
\tightlist
\item
  \textbf{Room protocol}: persistent SQLite chat; claim → red-team → formalize → relay-merge → independent cold audit. Relay cadence: every push to the working branch is merged to the mirror branch by the counterpart.
\item
  \textbf{Cold-audit rule}: ``verified'' means the file elaborates from source under the pinned toolchain (v4.31.0 + pinned mathlib) in an independent environment, unfiltered, with verbatim \texttt{\#print\ axioms} on the public theorems. Adopted after the stale-cache incident; upgraded after the rg-mask incident.
\item
  \textbf{Certificate factory}: exact DIMACS emitters with input hashes; kissat/cadical portfolios; DRAT/LRAT replay; compressed artifacts plus a manifest on durable storage; and resumable per-verdict queues. A SAT result becomes mathematical evidence only after a semantic bridge proves that every graph in the stated stratum satisfies the exact checked CNF. One recorded order-49 decomposition uses the checked-grid interface \texttt{orderFortyNineSmallHigh\_unsat\_of\_checkedCubeGrid}: 406 jobs comprising seven 7-by-8 positive-cube grids and fourteen negative-cover checks (room msgs 31965 and 31971). This historical decomposition is not a current pending-job count; the closure inventory must reconcile it with later covers and completed artifacts.
\item
  \textbf{Census tooling}: exhaustive B\&B sweeps over partition/atom spaces with every constraint tied to a named Lean lemma (loads, budgets, balance integrality, the equal-LCM law, oriented-cover kernels).
\end{itemize}

\section{Results and evidence map}\label{results-and-evidence-map}

The paper has two different logical endpoints and keeps them separate. Result A is a finite computational claim whose H1 verdict-only census is complete at this revision. It does not carry the status of an unconditional Lean theorem. Theorem B is a uniform conditional result: the negation of Erdős 85 follows from A-REG, while A-REG itself remains open. The implications from an unbounded family of plane-order drops to the negation of Erdős 85 are proved as \texttt{erdos85Negation\_iff\_not\_question}, \texttt{PlaneOrderDropWitness.strict\_drop}, and \texttt{not\_erdos85Question\_of\_cofinalPlaneOrderDropFamily}. On the binary branch, the existence jaw, tight-core reduction, and even-order regularity are proved by \texttt{Polarity.c4FreeMinDegreeWitness\_even\_delete\_absolute\_nucleus}, \texttt{binarySquareOrderTightCoreExclusion\_iff}, and \texttt{squareOrder\_regular\_of\_even}. The checked capstone is \texttt{not\_erdos85Question\_of\_binarySquareRegularExclusion}.

The strongest unconditional uniform reduction beneath A-REG is already substantial. The defect operator satisfies \texttt{A²\ =\ (q−1)I\ +\ J\ −\ D} and commutes with \texttt{A} (\texttt{adjMatrix\_sq\_eq\_sub\_secondOrderDefect\_of\_regular} and \texttt{adjMatrix\_comm\_secondOrderDefect\_of\_regular}). Its components have orders \texttt{q\ m\_c}, with \texttt{Σm\_c=q}; unit parts are impossible; and no component is bipartite when \texttt{4\ \textbar{}\ q} (\texttt{binarySquare\_regular\_exists\_defectComponent\_partition}, \texttt{binarySquare\_regular\_no\_sizeQ\_defectComponent\_of\_even}, and \texttt{binarySquare\_regular\_no\_bipartite\_defectComponent}). What remains is the all-non-bipartite connected-or-mixed node A-REG-NONBIP; the post-inverse divergence found no surviving terminal (room msgs 31962--31964).

\subsection{The computational 48-to-49 campaign}\label{the-computational-48-to-49-campaign}

Result A combines two checked witness ingredients---a 48-vertex degree-seven \texttt{C₄}-free witness and a 49-vertex degree-six witness---with heterogeneous evidence against a 49-vertex minimum-degree-seven witness. The graph-to-CNF consumer \texttt{not\_c4FreeMinDegreeWitness\_fortyNine\_seven\_of\_smallHighCubeBaseUnsat} assembles the one-, three-, five-, and seven-high exclusions. The finite-drop core would yield the exact thresholds and strict inequality if the nonexistence hypothesis were proved. The certificate universe, its structural decomposition, and the dependency-cone audit describe what formal verification would require; they are not a third headline result.

The order-48 existence input deserves separate notice. An independent NetworkX check over archived graph6 artifacts finds the checked extremal witness non-isomorphic to the Afzaly--McKay record, providing a second realization at the same extremal parameters. This comparison is reproducible computation rather than a Lean theorem; its script and graph6 inputs must be included in the release artifact and cited separately from the generated Lean endpoint. They are archived in \texttt{sat49/verify\_boza48\_nonisomorphism.py} and \texttt{sat49/data/}.

The campaign's operational scale explains the hybrid proof architecture. The authoritative H1 capacity universe contains 13,351 rows, while the higher strata decompose into checked cover and cube interfaces. SAT search produces LRAT evidence; Lean fixes the graph semantics, inventory bijections, and final composition. Receipts are resumability bookkeeping. The proof trust root is the kernel replay, clean cold integration build, and literal public-theorem axiom audit required by mandate 1318.

\subsection{The 63-to-64 campaign: useful finite evidence, still open}\label{the-63-to-64-campaign-useful-finite-evidence-still-open}

At \texttt{q=8}, the possible defect-component partitions are \texttt{{[}2,2,2,2{]}}, \texttt{{[}3,3,2{]}}, \texttt{{[}4,2,2{]}}, \texttt{{[}4,4{]}}, \texttt{{[}5,3{]}}, \texttt{{[}6,2{]}}, and \texttt{{[}8{]}}. They are not all excluded. The size-two \texttt{μ=3} sector is closed on honest regular hypotheses by \texttt{orderSixtyFour\_regular\_sizeTwoEigenline\_false}. The complete negative signed-joint size-two subtree, including the corrected \texttt{μ=-3,(0,5)} endpoint, is closed by \texttt{orderSixtyFour\_regular\_sizeTwo\_signedJoint\_false\_of\_connected}. But \texttt{{[}3,3,2{]}}, \texttt{{[}4,2,2{]}}, and \texttt{{[}6,2{]}} retain non-bipartite cases; \texttt{{[}4,4{]}} and \texttt{{[}5,3{]}} have only partial owner-nullity information; and the connected \texttt{{[}8{]}} case remains a gap. The eleven \texttt{{[}2,2,2,2{]}} assembly targets are external UNSAT verdicts without certificates. Therefore \texttt{63→64} is \textbf{not a decided drop}, and none of these order-64 enumerations proves the uniform A-REG statement (outline v2.64, §A.5.2).

\subsection{Plane-order censuses and the literature boundary}\label{plane-order-censuses-and-the-literature-boundary}

The order-80 search for a 9-regular \texttt{C₄}-free graph, which would supply the \texttt{q=9} existence jaw, is unresolved. All 20 authorized q9 solver attempts ended UNKNOWN under their caps; the independently checked positive controls returned SAT, but neither is an order-80 witness. The same campaign also left the order-78 existence cases open. Its 48-hour ledger is frozen in \texttt{q9\_solver\_controls/Q9\_EXISTENCE\_DECISION\_20260911.md}. An UNKNOWN is not evidence of nonexistence.

The complete finite-group Cayley census examined 52 groups of order 80, 47 of order 120, and 57 of order 168 with inverse-closed connection sets of degree 9, 11, and 13 respectively. Every input returned UNSAT; the reviewed encoding and table audits cover all these groups (\texttt{cayley-census-q11-q13/CAYLEY\_CENSUS\_Q11\_Q13\_20260913.md}). This excludes the specified Cayley families, not arbitrary graphs at those orders. The same census found one order-48, degree-7 Cayley witness among 52 groups.

The deductions from \href{https://doi.org/10.1016/j.disc.2016.12.005}{Zhang--Chen--Cheng's polarity-graph result} and \href{https://arxiv.org/abs/2409.12770}{Boza's bounds} put the adjacent star-Ramsey values at \texttt{r(109)\ ∈\ \{120,121\}} and \texttt{r(155)\ ∈\ \{168,169\}}, where \texttt{r(s)=R(C₄,K\_\{1,s\})} (\texttt{cayley-census-q11-q13/LITERATURE\_CHECK\_20260915.md}). The upper choice in either pair is equivalent to existence of the relevant 11- or 13-regular graph on 120 or 168 vertices. The Cayley census cannot select the unrestricted value. The cited literature check found no exact determination of these two values within its stated search scope; this is not a claim that none exists.

\subsection{Certificate case study: the h305 80-owner error and 88-owner repair}\label{certificate-case-study-the-h305-80-owner-error-and-88-owner-repair}

The \texttt{μ=-3,(0,5)} endpoint is a compact example of why hypothesis fidelity matters more than an UNSAT line. The first encoding silently reused the h114 shore table: eight fixed owners per shore and 80 candidate owners. The actual h305 shore modes contain antipodal offsets as well, giving twelve fixed edges per shore and an 88-owner universe. Consequently the original 80-owner UNSAT result excluded only a strengthened, wrong problem (room msgs 31664, 31674, and 31697; outline v2.64 change entry 2.61).

The repair rebuilt the whole chain around the honest table. Six canonical 88-owner CNFs were emitted independently and agreed byte-for-byte; all six were UNSAT, and their LRAT payloads are checked by the six \texttt{h305Owner88*\_check} theorems in \texttt{Erdos85MuNegThreeZeroFiveCorrectOwnerCertificate}. The graph semantics culminate in \texttt{muNegThreeZeroFiveCorrect\_graph\_false\_of\_exterior}, the source/transport endpoint in \texttt{false\_of\_h305\_source\_or\_transported}, and the callback-free order-64 consumer above. The six checked payloads contribute exactly six disclosed theorem-specific compiler-generated \texttt{native\_decide} axioms to this corrected endpoint; all structural bridge modules use standard axioms. The result is genuinely \texttt{PROVEN-AT-64\ CERT}, but only for that endpoint---not for the full order-64 nonexistence theorem (room msgs 31845--31846 and outline v2.64 entry 2.61).

\section{8. Interpretation: a computational drop, an honest infinite frontier, and the price of certainty}\label{interpretation-a-computational-drop-an-honest-infinite-frontier-and-the-price-of-certainty}

Result A is computational evidence, not a theorem, and this section is written so that it can be read standing alone. The lower sides of both values are explicit witnesses elaborated in Lean through \texttt{native\_decide}, which adds six named trust axioms beyond Lean's standard three; only Theorem B is standard-axiom-only in the cold-rebuild \texttt{\#print\ axioms} audit (\texttt{AXIOM\_AUDIT\_COLD\_20260927/}). The upper side at order 49 is a case split into the strata H1, H3, H5 and H7. Three of the four are closed at the level of a paper argument backed by reproducible, independently reviewed computation with banked receipts; the H1 rows were settled by a verdict-only census in which two independent SAT solvers had to return UNSAT under a declared cap. Rows that reached a cap were rerun at longer caps (4, then 12, then 24 hours per solver), and the single row that defeated every whole-instance cap was split by cube-and-conquer into 36 cubes that partition its search space exactly, each refuted by both solvers within an hour. No row is open: all 1,161 residual instances are UNSAT under two solvers. The honest statement is therefore that we have strong computational evidence that \texttt{f(48)\ =\ 8} and \texttt{f(49)\ =\ 7}, and therefore that the threshold drops between two adjacent orders. Erdős Problem 85 asks whether \texttt{f(n+1)\ ≥\ f(n)} holds for all large \texttt{n}. One drop at 48-to-49 is a data point against monotonicity at small order; it is compatible with either answer to the problem as posed, which a finite computation cannot settle, and we make no claim about the problem itself.

We chose to stop at belief and publish the price of certainty. A bank of 12,019 historical ready certificate inputs for the H1 rows is inventoried (objects and metadata, not yet validated payload by payload); replaying it through the Lean kernel is a priced, bounded task (about \$1,276 of spot infrastructure and 17 ideal allocated days on 16 hosts, plus about \$182 of retrieval), and the remaining rows have no credible fixed price because they were selected by difficulty. The certificate bank can be released as a requester-pays object store with an exact manifest, so that anyone who wants Result A promoted to a kernel-checked theorem can pay for exactly that and nothing else. The trust boundary is stated rather than hidden: two solvers agreeing under a cap is evidence about solver behaviour, not a proof; the H7 source enumeration rests on an audited program rather than an independent replay; and the archived H1 certificates were checked by \texttt{drat-trim} at production time, not by a current kernel replay.

Theorem B identifies what would turn a single drop into an infinite family. It is deliberately stated as a reduction from the one proposition A-REG. The evidence is mixed. Even-characteristic polarity graphs supply the cofinal existence jaw, and the binary-square defect calculus removes unit and bipartite components. Against that, the Lean-checked orders 15 and 16 and the published order-35/36 table entries show no drop, the \texttt{q\ =\ 4} analogue is false, and every attempted generic terminal for the remaining non-bipartite completion problem has failed or exposed a weaker hypothesis. A-REG is the live mathematical frontier, not a conclusion licensed by the finite data. The operator's plane-order interpretation is the principal rival: special orders may organize both the constructions and the obstructions without the binary regular-exclusion pattern persisting uniformly. Publishing the twin result makes that disagreement useful: the checked reduction states exactly what a proof of the negative answer must supply, the negative map states which plausible shortcuts do not supply it, and Result A provides a concrete calibration point for future theory.

This section was reconciled on 2026-09-28 with the three artifacts it names and nothing else. The receipt-derived census: the reviewed Phase B summarizer over 1,413 run directories and the Mac pilot (\texttt{phase\_b\_h1\_census\_20260927/h1-census-table.json}: 1,160 two-solver UNSAT rows, 96 historical rows, no disagreement) and the reviewed 1,288-slot gap auditor (\texttt{h1-gap1288-audit.json}: 1,191 fresh two-solver UNSAT slots, 96 inherited, one cube-partitioned slot reported separately by \texttt{cube-h1\_81494a6ef36d3ec9/tree-check.json}); the open list is empty. The 12,019 certificate rows enter only as the inventory that prices replay; they are not re-joined here. The literal \texttt{\#print\ axioms} output from a cold rebuild for the witnesses, the conditional finite-drop core and Theorem B (\texttt{AXIOM\_AUDIT\_COLD\_20260927/}). The release manifest of the certificate bank is not part of this revision because no requester-pays copy has been published; if one is, its manifest must be added here. No operational receipt, fleet count or verdict summary substitutes for those artifacts.

\bibliographystyle{plainnat}
\bibliography{refs}

\end{document}
